\documentclass[UTF8,a4paper,11pt,fontset=fandol]{ctexart}
\newcommand{\DocTitle}{实验01\_实验手册}
\newcommand{\DocSubject}{2026年秋季课程B：第一次实验，Python环境与基础入门}
\newcommand{\DocShortTitle}{实验一：环境与基础入门}
\newcommand{\DocType}{学生实验手册}
% Shared style for integrated experiment handouts.
\usepackage[left=2.15cm,right=2.15cm,top=2.25cm,bottom=2.15cm,headsep=0.7cm]{geometry}
\usepackage{amsmath,amssymb}
\usepackage{booktabs,array,tabularx,longtable,multirow}
\usepackage{enumitem}
\usepackage{xcolor}
\usepackage{graphicx}
\usepackage{tikz}
\usetikzlibrary{arrows.meta,positioning,calc,shapes.geometric}
\usepackage{tcolorbox}
\usepackage{listings}
\usepackage{hyperref}
\usepackage{fancyhdr}
\usepackage{chngcntr}

\definecolor{Ink}{HTML}{17233C}
\definecolor{Navy}{HTML}{173B68}
\definecolor{Blue}{HTML}{2878B5}
\definecolor{Cyan}{HTML}{2A9D8F}
\definecolor{Amber}{HTML}{E9A23B}
\definecolor{Red}{HTML}{C84B4B}
\definecolor{Paper}{HTML}{F6F8FB}
\definecolor{PaleBlue}{HTML}{EAF2FA}
\definecolor{PaleCyan}{HTML}{EAF7F4}
\definecolor{PaleAmber}{HTML}{FFF5E4}
\definecolor{Rule}{HTML}{D6DEE8}
\definecolor{Muted}{HTML}{5B677A}

\hypersetup{
  colorlinks=true,
  linkcolor=Navy,
  urlcolor=Blue,
  citecolor=Cyan,
  pdfauthor={《人工智能数学原理与算法》课程组},
  pdftitle={\DocTitle},
  pdfsubject={\DocSubject}
}

\ctexset{
  section={format=\Large\bfseries\color{Navy},beforeskip=1.4ex plus .2ex,afterskip=.8ex},
  subsection={format=\large\bfseries\color{Ink},beforeskip=1.2ex plus .2ex,afterskip=.55ex},
  subsubsection={format=\normalsize\bfseries\color{Blue},beforeskip=1ex,afterskip=.35ex}
}

\setlength{\parindent}{2em}
\setlength{\parskip}{0.22em}
\setlength{\emergencystretch}{2em}
\linespread{1.14}
\setlist[itemize]{leftmargin=2em,itemsep=.25em,topsep=.35em}
\setlist[enumerate]{leftmargin=2.2em,itemsep=.28em,topsep=.35em}
\renewcommand{\arraystretch}{1.30}
\newcolumntype{Y}{>{\raggedright\arraybackslash}X}
\newcolumntype{C}[1]{>{\centering\arraybackslash}p{#1}}
\newcolumntype{L}[1]{>{\raggedright\arraybackslash}p{#1}}

\newcommand{\R}{\mathbb{R}}
\newcommand{\E}{\mathbb{E}}
\newcommand{\MSE}{\operatorname{MSE}}
\newcommand{\RMSE}{\operatorname{RMSE}}
\newcommand{\argmin}{\operatorname*{arg\,min}}
\newcommand{\vect}[1]{\boldsymbol{#1}}
\newcommand{\mat}[1]{\mathbf{#1}}
\newcommand{\term}[1]{\textcolor{Blue}{\textbf{#1}}}
\newcommand{\code}[1]{\texttt{#1}}
\numberwithin{equation}{subsection}

\tcbset{boxrule=.7pt,arc=1.6mm,left=2.5mm,right=2.5mm,top=1.8mm,bottom=1.8mm,before skip=7pt,after skip=7pt}
\newtcolorbox{keybox}[1][核心结论]{colback=PaleBlue,colframe=Blue,title={#1},fonttitle=\bfseries,coltitle=Ink}
\newtcolorbox{examplebox}[1][贯穿例题]{colback=PaleCyan,colframe=Cyan,title={#1},fonttitle=\bfseries,coltitle=Ink}
\newtcolorbox{questionbox}[1][检查点]{colback=PaleAmber,colframe=Amber,title={#1},fonttitle=\bfseries,coltitle=Ink}
\newtcolorbox{pitfallbox}[1][易错提醒]{colback=red!3,colframe=Red,title={#1},fonttitle=\bfseries,coltitle=Ink}
\newtcolorbox{teacherbox}[1][教师提示]{colback=Paper,colframe=Rule,title={#1},fonttitle=\bfseries,coltitle=Muted}
\newtcolorbox{proofbox}[1][推导与证明]{colback=Blue!3,colframe=Navy,title={#1},fonttitle=\bfseries,coltitle=Ink}

\lstdefinestyle{python}{
  language=Python,
  basicstyle=\ttfamily\small,
  keywordstyle=\color{Blue}\bfseries,
  commentstyle=\color{Cyan!70!black},
  stringstyle=\color{Red!80!black},
  showstringspaces=false,
  columns=fullflexible,
  keepspaces=true,
  frame=single,
  rulecolor=\color{Rule},
  backgroundcolor=\color{Paper},
  breaklines=true,
  numbers=left,
  numberstyle=\tiny\color{Muted},
  xleftmargin=1.5em,
  framexleftmargin=1.2em,
  aboveskip=.7em,
  belowskip=.7em
}
\lstset{style=python}

\pagestyle{fancy}
\fancyhf{}
\fancyhead[L]{\small\color{Muted}人工智能数学原理与算法}
\fancyhead[R]{\small\color{Muted}\DocShortTitle}
\fancyfoot[L]{\small\color{Muted}\DocType}
\fancyfoot[R]{\small\color{Muted}\thepage}
\renewcommand{\headrulewidth}{0pt}
\renewcommand{\footrulewidth}{0pt}

\newcommand{\makeexperimenttitle}[4]{
  \hypersetup{pageanchor=false}
  \begin{titlepage}\thispagestyle{empty}
  \begin{tikzpicture}[remember picture,overlay]
    \fill[Ink] (current page.north west) rectangle ([yshift=-9.4cm]current page.north east);
    \fill[Blue] ([yshift=-9.4cm]current page.north west) rectangle ([yshift=-9.68cm]current page.north east);
    \fill[Cyan] ([xshift=2.15cm,yshift=-1.55cm]current page.north west) rectangle ([xshift=2.27cm,yshift=-8.55cm]current page.north west);
    \fill[Paper] ([xshift=2.15cm,yshift=2.0cm]current page.south west) rectangle ([xshift=-2.15cm,yshift=6.15cm]current page.south east);
  \end{tikzpicture}
  \vspace*{1.45cm}
  {\color{white}\Large 人工智能数学原理与算法\par}\vspace{.55cm}
  {\color{white}\fontsize{31}{39}\selectfont\bfseries #1\par}\vspace{.35cm}
  {\color{white}\fontsize{21}{28}\selectfont #2\par}
  \vfill
  \begin{tcolorbox}[colback=white,colframe=white,boxrule=0pt,arc=0mm,left=7mm,right=7mm,top=6mm,bottom=6mm,width=\textwidth]
    \begin{tabularx}{\linewidth}{L{2.6cm}Y}
      \textcolor{Muted}{材料类型} & \DocType\\
      \textcolor{Muted}{实验安排} & #3\\
      \textcolor{Muted}{贯穿案例} & #4\\
      \textcolor{Muted}{适用对象} & 已完成机器学习导论、具备高中数学与基础计算机操作能力的本科生\\
    \end{tabularx}
  \end{tcolorbox}
  \vspace{.7cm}{\color{Muted}\small 版本：2026年秋季学期整合稿\hfill 源文件与 PDF 同目录交付\par}
  \end{titlepage}
  \hypersetup{pageanchor=true}
  \pagenumbering{roman}\tableofcontents\clearpage\pagenumbering{arabic}
}

\usepackage{calc}
\usepackage{etoolbox}
\usepackage{xurl}
\usepackage{bookmark}
\setmonofont[BoldFont=DejaVuSansMono-Bold.ttf,ItalicFont=DejaVuSansMono-Oblique.ttf,BoldItalicFont=DejaVuSansMono-BoldOblique.ttf]{DejaVuSansMono.ttf}
\setCJKmonofont[AutoFakeBold=2]{FandolFang-Regular.otf}
\newcounter{none}
\providecommand{\tightlist}{\setlength{\itemsep}{0pt}\setlength{\parskip}{0pt}}
\providecommand{\passthrough}[1]{#1}
\providecommand{\pandocbounded}[1]{#1}
\setcounter{tocdepth}{1}
\setcounter{secnumdepth}{2}
\setlength{\headheight}{15pt}
\lstset{basicstyle=\ttfamily\small,numbers=none,xleftmargin=0pt,framexleftmargin=0pt,breakatwhitespace=false}
\AtBeginEnvironment{longtable}{\small\setlength{\tabcolsep}{4pt}}
\hypersetup{hypertexnames=false,pdfauthor={钟志诚（主讲）；曾有成、史睿铭、胡泷（整理）}}
\begin{document}
\raggedright
\setlength{\parindent}{2em}
\hypersetup{pageanchor=false}
\begin{titlepage}
\thispagestyle{empty}
\vspace*{1.5cm}
{\Large\color{Navy}人工智能数学原理与算法 B\par}
\vspace{1cm}
{\huge\bfseries\color{Ink}第一次实验\par}
\vspace{.45cm}
{\LARGE\bfseries\color{Ink}Python 环境与基础入门\par}
\vspace{.6cm}
{\Large 学生实验手册\par}
\vspace{1cm}
{\large 2026 年秋季学期\par}
\vspace{.8cm}
\begin{tabular}{@{}ll}
主讲 & 钟志诚\\[.35em]
整理 & 曾有成、史睿铭、胡泷\\[.35em]
单位 & 中国科学技术大学\\
     & 人工智能与数据科学学院
\end{tabular}
\vspace{1cm}\par
建议课前完成环境安装与首次运行。\\
日期、地点与提交平台以课程通知为准。\par
\vfill
\noindent\rule{\linewidth}{.5pt}\par
微型书店：对象、继承、数据管理与销售统计。\\
鸢尾花：补全模型调用，完成训练与结果验证。\\
基础练习 Notebook 配合教学，不要求提交。\par
\end{titlepage}
\hypersetup{pageanchor=true}
\pagenumbering{roman}
\tableofcontents
\clearpage
\pagenumbering{arabic}
\section{实验内容与课前准备}\label{ux5b9eux9a8cux5185ux5bb9ux4e0eux8bfeux524dux51c6ux5907}

在本次实验中，我们将使用 Python 完成几项书目处理任务，并运行一个简单的分类示例。书目任务帮助你熟悉数据组织、条件判断、循环和函数；后面的数组、对象和数据探索练习，则为继续使用机器学习库作准备。相关概念可以结合\href{实验01_实验讲义.md}{实验讲义}阅读。

\textbf{建议在课前完成环境安装与首次运行，尤其是参加线下实验课的同学。} 集中授课时，网络质量可能影响安装包和依赖的下载速度与稳定性。请尽量在网络较好的时候完成本手册第 2--4 节，提前发现问题；课堂上将核对运行情况，并集中处理仍未解决的问题。

整个准备过程包括下载 Miniforge 安装包、创建 Python 环境、安装实验所需的库和运行程序。前三项都可能需要下载文件，因此只保存安装包还不算准备完成。程序和依赖准备好后，本次的书目与鸢尾花示例无需再下载数据。

本次实验无需显卡。已有 Anaconda、Miniconda 或 Miniforge 的同学可以从第 2.1 节检查现有安装；使用学校预配置电脑时，先按学校提供的方式进入环境，再执行第 3.3 节的检查。

如果课前未能完成安装，请带上已经下载的材料，并记录执行的命令和报错。尚未解决的问题可以在课堂继续处理；环境准备成功的同学可先运行和修改基础例子。

\clearpage

\section{安装或检查 conda}\label{ux5b89ux88c5ux6216ux68c0ux67e5-conda}

Python 解释器负责运行程序，NumPy、scikit-learn 等库提供数值计算和模型接口。conda 用于管理 Python 环境，让本课程使用的解释器和依赖与其他项目分开。Miniforge 是包含 conda 的安装发行版；装好后，还需要创建课程环境并安装实验依赖。

\subsection{已经安装过的同学}\label{ux5df2ux7ecfux5b89ux88c5ux8fc7ux7684ux540cux5b66}

Windows 用户打开已有的 \textbf{Miniforge Prompt} 或 \textbf{Anaconda Prompt}；macOS/Linux 用户打开终端，执行：

\begin{lstlisting}[language={}]
conda --version
conda env list
\end{lstlisting}

若能看到 conda 的版本和环境列表，可以跳到第 3 节。若列表中已经有 \passthrough{\lstinline!ustc-ai-experiments-26fall-B!}，直接激活并检查依赖即可。

Windows 普通终端提示找不到 \passthrough{\lstinline!conda!}，不一定是没有安装。先在开始菜单中搜索上述 Prompt，再作判断。仅能执行 \passthrough{\lstinline!python!} 也不等于已安装 conda；不熟悉环境管理时，可以按下文统一使用 Miniforge，不必先卸载原来的 Python。

\subsection{选择安装包}\label{ux9009ux62e9ux5b89ux88c5ux5305}

从 \href{https://github.com/conda-forge/miniforge\#requirements-and-installers}{Miniforge 官方页面}选择与你的操作系统和处理器匹配的安装包。GitHub 下载缓慢时，可使用\href{https://mirrors.ustc.edu.cn/github-release/conda-forge/miniforge/LatestRelease/}{科大 Miniforge 镜像目录}。下载页中的安装器可能带版本号，认准系统和架构即可。

\begin{longtable}[]{@{}ll@{}}
\toprule\noalign{}
电脑 & 常见安装包名称 \\
\midrule\noalign{}
\endhead
\bottomrule\noalign{}
\endlastfoot
Windows，Intel/AMD 64 位处理器 & \passthrough{\lstinline!Miniforge3-Windows-x86\_64.exe!} \\
macOS，Apple 芯片 & \passthrough{\lstinline!Miniforge3-MacOSX-arm64.sh!} \\
macOS，Intel 处理器 & \passthrough{\lstinline!Miniforge3-MacOSX-x86\_64.sh!} \\
Linux，Intel/AMD 64 位处理器 & \passthrough{\lstinline!Miniforge3-Linux-x86\_64.sh!} \\
Linux，ARM 64 位处理器 & \passthrough{\lstinline!Miniforge3-Linux-aarch64.sh!} \\
\end{longtable}

Windows 可在``设置 → 系统 → 系统信息（关于）''中查看系统类型；Mac 可在苹果菜单的``关于本机''中查看芯片或处理器。Linux 可执行 \passthrough{\lstinline!uname -m!} 查看架构。Windows ARM 电脑或较旧系统请先核对官方支持范围，不能只看文件名中有 ``Windows'' 就直接选用。

以下步骤依据 \href{https://github.com/conda-forge/miniforge\#install}{Miniforge 安装说明}。

\subsection{Windows 安装}\label{windows-ux5b89ux88c5}

\begin{enumerate}
\def\labelenumi{\arabic{enumi}.}
\tightlist
\item
  双击下载的 \passthrough{\lstinline!.exe!} 文件，阅读许可说明并继续。安装范围通常选择当前用户（\passthrough{\lstinline!Just Me!}）。
\item
  选择自己有写入权限、没有空格和中文等特殊字符的安装路径，例如 \passthrough{\lstinline!D:\\Tools\\miniforge3!}。没有 D 盘时使用其他可写位置，不要照抄不存在的盘符。
\item
  保留创建开始菜单快捷方式的选项。通过 Miniforge Prompt 使用时，不需要勾选 \passthrough{\lstinline!Add Miniforge3 to my PATH environment variable!}；也无需把它注册成系统默认 Python。
\item
  安装完成后，从开始菜单打开 \textbf{Miniforge Prompt}。窗口中通常能看到 \passthrough{\lstinline!(base)!}，表示当前使用 Miniforge 的基础环境。
\item
  执行第 2.1 节的两条命令，确认能显示版本与环境列表。
\end{enumerate}

后面的终端命令都可以在这个窗口执行。每次重新打开窗口后，需要重新激活课程环境。暂时无需为 PowerShell 或编辑器另外配置 conda。

\subsection{macOS / Linux 安装}\label{macos-linux-ux5b89ux88c5}

下面使用 \passthrough{\lstinline!.sh!} 安装器。打开终端，进入下载目录；若浏览器把文件放在别处，请换成实际目录：

\begin{lstlisting}[language={}]
cd ~/Downloads
ls
\end{lstlisting}

找到下载的文件后，用 \passthrough{\lstinline!bash!} 执行。以 Apple 芯片的 Mac 为例：

\begin{lstlisting}[language={}]
bash Miniforge3-MacOSX-arm64.sh
\end{lstlisting}

Intel Mac 或 Linux 用户将文件名替换为自己的安装包名称。按提示阅读并确认许可，选择安装位置；个人电脑通常可使用默认的 \passthrough{\lstinline!\~/miniforge3!}。安装器询问是否初始化 conda 时选择 \passthrough{\lstinline!yes!}，让新开的终端能够使用 conda。

安装结束后关闭并重新打开终端，执行第 2.1 节的检查。如果安装时没有初始化，可先在当前终端加载 conda，路径对应默认安装位置：

\begin{lstlisting}[language={}]
source ~/miniforge3/etc/profile.d/conda.sh
\end{lstlisting}

这条命令只让当前终端能够使用 conda；以后新开的终端也需要加载。若想固定初始化设置，可再按官方说明处理，或把安装位置和终端类型提供给助教。

\clearpage

\section{创建课程环境并安装依赖}\label{ux521bux5efaux8bfeux7a0bux73afux5883ux5e76ux5b89ux88c5ux4f9dux8d56}

以下命令在终端中逐行执行。等待上一条结束、重新出现命令提示符后，再输入下一条。若窗口中只有 Python 的 \passthrough{\lstinline!>>>!}，先输入 \passthrough{\lstinline!exit()!} 返回终端。

\subsection{创建与激活}\label{ux521bux5efaux4e0eux6fc0ux6d3b}

先用 \passthrough{\lstinline!conda env list!} 检查是否已有课程环境。没有时执行：

\begin{lstlisting}[language={}]
conda create -n ustc-ai-experiments-26fall-B --override-channels -c conda-forge python=3.11 pip
\end{lstlisting}

\passthrough{\lstinline!-n!} 后面是环境名；\passthrough{\lstinline!python=3.11!} 指定课程使用的 Python 版本。这里明确使用 \passthrough{\lstinline!conda-forge!} 下载 Python 和 pip，\passthrough{\lstinline!--override-channels!} 使本次创建只使用指定频道，不改动你保存的全局频道设置。

conda 会列出准备安装的软件，询问 \passthrough{\lstinline!Proceed ([y]/n)?!} 时输入 \passthrough{\lstinline!y!} 并回车。安装完成后激活环境：

\begin{lstlisting}[language={}]
conda activate ustc-ai-experiments-26fall-B
\end{lstlisting}

终端提示符前通常会出现课程环境名。已有同名环境时只需激活，不要再次创建。Miniforge 基础环境的 Python 版本可能不同，以激活后的检查结果为准。

\subsection{安装实验依赖}\label{ux5b89ux88c5ux5b9eux9a8cux4f9dux8d56}

在已激活的课程环境中执行：

\begin{lstlisting}[language={}]
python -m pip install numpy jupyterlab scikit-learn matplotlib
\end{lstlisting}

\begin{longtable}[]{@{}ll@{}}
\toprule\noalign{}
库 & 本次用途 \\
\midrule\noalign{}
\endhead
\bottomrule\noalign{}
\endlastfoot
NumPy & 数组与数值计算 \\
JupyterLab & 在浏览器中编辑和运行 Notebook \\
scikit-learn & 内置鸢尾花数据、特征处理与分类模型 \\
Matplotlib & 绘制散点图和直方图 \\
\end{longtable}

\passthrough{\lstinline!python -m pip!} 使用当前 Python 对应的安装工具，有助于把库安装到刚才激活的环境中。安装名是 \passthrough{\lstinline!scikit-learn!}，代码中使用的导入名则是 \passthrough{\lstinline!sklearn!}。

安装结束后应回到命令提示符，且没有安装失败的 \passthrough{\lstinline!ERROR!}。看到 \passthrough{\lstinline!Requirement already satisfied!} 表示对应依赖已满足，可以继续检查。JupyterLab 的安装方式也可查阅\href{https://jupyterlab.readthedocs.io/en/stable/getting_started/installation.html}{官方说明}。

\subsection{检查解释器和依赖}\label{ux68c0ux67e5ux89e3ux91caux5668ux548cux4f9dux8d56}

执行以下命令：

\begin{lstlisting}[language={}]
python --version
python -c "import sys; print(sys.executable)"
python -c "import numpy as np; print(np.__version__); print(np.array([1, 2, 3]) + 1)"
python -c "import sklearn; print(sklearn.__version__)"
python -c "import matplotlib; print(matplotlib.__version__)"
python -m jupyterlab --version
\end{lstlisting}

依次检查：

\begin{itemize}
\tightlist
\item
  Python 版本应为 \passthrough{\lstinline!3.11.x!}；学校统一提供的其他已测试版本，以课程配置为准。
\item
  解释器路径应指向课程环境。例如 Windows 路径中通常包含 \passthrough{\lstinline!envs\\ustc-ai-experiments-26fall-B\\python.exe!}，而不是另一个项目的 Python。
\item
  NumPy 命令应输出版本号和 \passthrough{\lstinline![2 3 4]!}。
\item
  其余命令应能显示相应版本，不出现 \passthrough{\lstinline!ModuleNotFoundError!}。
\end{itemize}

具体库版本可能不同，版本号本身不是评分点。请记下解释器路径，后面可以用它核对编辑器和 Notebook 是否使用同一个环境。

\subsection{下载缓慢或中断时}\label{ux4e0bux8f7dux7f13ux6162ux6216ux4e2dux65adux65f6}

先区分卡在哪个环节，再选择对应办法。换源不能解决所有网络问题，也不会省去下载本身。

\textbf{下载 Miniforge 安装包较慢：} 可使用第 2.2 节的科大镜像目录，仍需选择相同系统和架构的安装器。

\textbf{\passthrough{\lstinline!conda create!} 下载缓慢或连接超时：} 在原命令已经结束、且课程环境尚未成功创建时，可改用科大的 conda-forge 镜像重试：

\begin{lstlisting}[language={}]
conda create -n ustc-ai-experiments-26fall-B --override-channels -c https://mirrors.ustc.edu.cn/anaconda/cloud/conda-forge python=3.11 pip
\end{lstlisting}

两条创建命令选用一条即可。若提示同名环境或路径已存在，先执行 \passthrough{\lstinline!conda env list!}，尝试激活后检查 Python；仍有问题时保留报错，请助教协助检查，不必先删除整个安装目录。

\textbf{\passthrough{\lstinline!pip install!} 下载缓慢或超时：} 确认课程环境已激活，再执行：

\begin{lstlisting}[language={}]
python -m pip install -i https://mirrors.ustc.edu.cn/pypi/simple numpy jupyterlab scikit-learn matplotlib
\end{lstlisting}

这些备用命令仅为本次操作指定下载地址。镜像用法见\href{https://mirrors.ustc.edu.cn/help/anaconda.html}{科大 conda 镜像说明}和\href{https://mirrors.ustc.edu.cn/help/pypi.html}{PyPI 镜像说明}。

如果出现证书错误、访问被拒绝或连续下载失败，请先完成校园网登录、核对系统时间，记录报错和网络环境后求助。无需通过关闭证书校验来完成安装。

\clearpage

\section{下载材料并首次运行}\label{ux4e0bux8f7dux6750ux6599ux5e76ux9996ux6b21ux8fd0ux884c}

\subsection{解压并找到练习与项目}\label{ux89e3ux538bux5e76ux627eux5230ux7ec3ux4e60ux4e0eux9879ux76ee}

从\href{../downloads/lab01-python-basics-26fall-B.zip}{学生材料包}下载 ZIP，完整解压到你准备保存实验的目录。不要直接在压缩包预览中运行程序。找到 \passthrough{\lstinline!实验01\_Python环境与基础入门!} 下的 \passthrough{\lstinline!code!} 文件夹：

\begin{lstlisting}[language={}]
code/
├── README.md
├── practice.ipynb
├── books.txt
├── bookstore/
│   ├── README.md
│   ├── book_tasks.py
│   ├── book_support.py
│   ├── check_book_tasks.py
│   ├── books.csv
│   └── sales.csv
└── iris/
    ├── README.md
    ├── iris_classification.py
    └── check_iris.py
\end{lstlisting}

\passthrough{\lstinline!practice.ipynb!} 用于 Python 基础练习，不要求提交。两个微项目各自包含说明和代码，书店的数据也保存在项目内。\passthrough{\lstinline!books.txt!} 仅供讲义第 5 章演示文本读取，与书店项目的 CSV 输入无关。

\subsection{在终端中进入 code 文件夹}\label{ux5728ux7ec8ux7aefux4e2dux8fdbux5165-code-ux6587ux4ef6ux5939}

打开终端并激活课程环境：

\begin{lstlisting}[language={}]
conda activate ustc-ai-experiments-26fall-B
\end{lstlisting}

\passthrough{\lstinline!cd!} 用来切换当前工作目录。先进入 \passthrough{\lstinline!code!}，即可打开基础练习或访问两个项目子目录。下面的路径只是例子，需替换为你实际解压得到的路径。

\textbf{Windows Miniforge Prompt / 命令提示符：}

\begin{lstlisting}[language={}]
cd /d "D:\course\实验01_Python环境与基础入门\code"
dir
\end{lstlisting}

\passthrough{\lstinline!/d!} 允许同时切换盘符。如果你使用的是已经配置好 conda 的 PowerShell，则使用下面的写法：

\begin{lstlisting}[language={}]
cd "D:\course\实验01_Python环境与基础入门\code"
dir
\end{lstlisting}

\textbf{macOS / Linux 终端：}

\begin{lstlisting}[language={}]
cd "$HOME/course/实验01_Python环境与基础入门/code"
ls
\end{lstlisting}

列出的文件中应能看到 \passthrough{\lstinline!practice.ipynb!}、\passthrough{\lstinline!bookstore!} 和 \passthrough{\lstinline!iris!}。如果没有，请先检查当前位置。路径中有空格时需保留引号。

\clearpage

\subsection{先试运行一个 Python 脚本}\label{ux5148ux8bd5ux8fd0ux884cux4e00ux4e2a-python-ux811aux672c}

在 \passthrough{\lstinline!code!} 目录中新建 \passthrough{\lstinline!hello.py!}，用文本编辑器写入并保存：

\begin{lstlisting}[language=Python]
print("Hello, Python!")
\end{lstlisting}

在终端执行：

\begin{lstlisting}[language={}]
python hello.py
\end{lstlisting}

应看到 \passthrough{\lstinline"Hello, Python!"}。试着修改问候语，保存后重新运行，确认输出变化。这个文件是你自己创建的运行练习，不属于正式作业，也无需提交。Windows 用户注意查看扩展名，保存的文件应为 \passthrough{\lstinline!.py!}，而不是 \passthrough{\lstinline!.py.txt!}。

后续项目也用同样方式运行，只需指定项目脚本的路径，见第 5 节。若使用编辑器中的运行按钮，应选择与第 3.3 节一致的解释器；结果不一致时先用终端核对。

\subsection{运行基础练习 Notebook}\label{ux8fd0ux884cux57faux7840ux7ec3ux4e60-notebook}

在已激活环境、且位于 \passthrough{\lstinline!code!} 文件夹的终端中执行：

\begin{lstlisting}[language={}]
python -m jupyterlab
\end{lstlisting}

浏览器通常会自动打开 JupyterLab。若没有打开，可复制终端显示的本地地址到浏览器，保留地址中的登录参数；不要把带有令牌的完整地址发到公开群聊。JupyterLab 使用本机 Python 运行代码，打开本地界面不需要注册在线账户。参见\href{https://jupyterlab.readthedocs.io/en/stable/getting_started/starting.html}{启动说明}。

可以先打开 \href{code/practice.ipynb}{practice.ipynb}，运行准备单元和第 1 节，再按需要继续。它配合 Python 基础教学，不要求提交。若提示选择内核，选择当前环境提供的 Python 内核，常见名称为 \passthrough{\lstinline!Python 3 (ipykernel)!}。名称本身不足以确认环境，需运行检查代码：

\begin{lstlisting}[language=Python]
import sys
print(sys.executable)
\end{lstlisting}

输出应与第 3.3 节的路径一致。若不同，先检查启动 JupyterLab 前是否激活了课程环境，再从内核菜单选择对应内核；无法确定时，将两个路径提供给助教。

选中代码单元后按 \passthrough{\lstinline!Shift+Enter!}，运行并移动到下一格。第一次按从上往下的顺序执行；说明单元用于阅读，不执行 Python。修改函数定义后，先重新运行定义单元，再运行调用它的单元。不同 Notebook 不共享变量。

保存修改后，执行菜单中的``重启内核并运行所有单元''（Restart Kernel and Run All Cells），检查是否依赖了之前残留的变量。结束使用时先保存文件，再到启动 JupyterLab 的终端按 \passthrough{\lstinline!Ctrl+C!}，按提示确认停止服务。仅关闭浏览器标签页不一定会停止后台程序。

\subsection{课前检查到这里即可}\label{ux8bfeux524dux68c0ux67e5ux5230ux8fd9ux91ccux5373ux53ef}

请确认：

\begin{itemize}
\tightlist
\item
  能激活课程环境，知道当前使用的 Python 路径。
\item
  第 3.3 节的依赖检查通过。
\item
  能运行自己创建的 \passthrough{\lstinline!hello.py!} 和 \passthrough{\lstinline!practice.ipynb!}，修改问候语或第 1 节的变量后看到输出变化。
\end{itemize}

课前暂不要求完成两个编程任务。提前跑通程序是为了确认工具可用，后续再按项目说明完成任务。

\clearpage

\section{两个编程任务}\label{ux4e24ux4e2aux7f16ux7a0bux4efbux52a1}

本次作业包含两个小项目。每个项目都有独立说明，介绍问题、提供的材料、需要补全的功能和完成后的结果。建议先读项目说明，再打开代码搜索 \passthrough{\lstinline!LAB01\_TODO!}，完成一个位置就运行一次。

\begin{longtable}[]{@{}
  >{\raggedright\arraybackslash}p{(\linewidth - 4\tabcolsep) * \real{0.3333}}
  >{\raggedright\arraybackslash}p{(\linewidth - 4\tabcolsep) * \real{0.3333}}
  >{\raggedright\arraybackslash}p{(\linewidth - 4\tabcolsep) * \real{0.3333}}@{}}
\toprule\noalign{}
\begin{minipage}[b]{\linewidth}\raggedright
项目说明
\end{minipage} & \begin{minipage}[b]{\linewidth}\raggedright
作业文件
\end{minipage} & \begin{minipage}[b]{\linewidth}\raggedright
学生完成的部分
\end{minipage} \\
\midrule\noalign{}
\endhead
\bottomrule\noalign{}
\endlastfoot
\href{code/bookstore/README.md}{任务一：微型书店数据管理与统计} & \href{code/bookstore/book_tasks.py}{book\_tasks.py} & B1--B6 图书对象、更新、查询、销售解析与统计 \\
\href{code/iris/README.md}{任务二：鸢尾花数据观察与分类} & \href{code/iris/iris_classification.py}{iris\_classification.py} & I1 选择绘图特征、I2 调用训练、I3 调用预测 \\
\end{longtable}

书店任务使用对象、继承、条件与循环、文件记录和数组计算实现连续的数据管理功能。鸢尾花任务则练习连接现成工具：数据读取、绘图、模型创建与评价已经提供，项目说明会介绍调用所需的概念。

进入作业框架的 \passthrough{\lstinline!code!} 目录后运行：

\begin{lstlisting}[language={}]
python bookstore/book_tasks.py
python iris/iris_classification.py
\end{lstlisting}

I1 需要把两个 \passthrough{\lstinline!None!} 改成特征列编号；书店的 B1--B6 及鸢尾花 I2/I3 需要用实现替换 \passthrough{\lstinline!raise NotImplementedError(...)!} 占位语句。未完成时程序会提示对应位置；这类提示与环境安装失败不同。具体输入、输出和检查方法见各自项目说明。

\subsection{基础练习与补充阅读}\label{ux57faux7840ux7ec3ux4e60ux4e0eux8865ux5145ux9605ux8bfb}

\href{code/practice.ipynb}{practice.ipynb} 用短例子练习变量、容器、控制流、函数、对象与继承，以及模块和文件操作。它不要求提交，也不重复两个项目的功能补全。概念解释、C / Python 对照和选读示例保留在讲义中，按需查阅。

\section{建议的实现顺序}\label{ux5efaux8baeux7684ux5b9eux73b0ux987aux5e8f}

先完成环境检查，确认能运行基础练习，再按两个项目说明补全代码。

\begin{enumerate}
\def\labelenumi{\arabic{enumi}.}
\tightlist
\item
  先实现书店的图书对象、补货与改价，运行相应阶段的验证。完成对象初始化后，可以进入菜单查看书库。
\item
  完成查询和销售记录解析，观察公共读取父类怎样调用不同子类的方法。
\item
  完成销售记录匹配与数组计算，检查每种书和总体的销售额、毛利；再通过菜单修改售价、补货和保存，检查功能之间的联系。
\item
  完成鸢尾花项目的特征选择、训练和预测调用，观察图与评价结果。
\end{enumerate}

书店验证程序可以通过 \passthrough{\lstinline!--stage B3!} 等参数单独检查一个阶段，不需要等所有 TODO 都完成。遇到问题时，先确认正在处理的数据和对应函数，必要时参考讲义中的相关例子。

讲义中带 \passthrough{\lstinline!*!} 的内容供补充学习。基础练习、项目实现和补充阅读按自己的进度衔接，无需按固定时段完成。

\clearpage

\section{提交文件与提交前检查}\label{ux63d0ux4ea4ux6587ux4ef6ux4e0eux63d0ux4ea4ux524dux68c0ux67e5}

\subsection{提交一个 ZIP 压缩包}\label{ux63d0ux4ea4ux4e00ux4e2a-zip-ux538bux7f29ux5305}

将下面的文件放进一个文件夹，再压缩为 \textbf{\passthrough{\lstinline!学号\_姓名\_lab01.zip!}}。将``学号''和``姓名''替换成自己的信息，目录结构如下：

\begin{lstlisting}[language={}]
学号_姓名_lab01.zip
└── 学号_姓名_lab01/
    ├── code/
    │   ├── bookstore/
    │   │   ├── book_tasks.py
    │   │   ├── book_support.py
    │   │   ├── books.csv
    │   │   └── sales.csv
    │   └── iris/
    │       └── iris_classification.py
    └── report.pdf
\end{lstlisting}

\begin{longtable}[]{@{}
  >{\raggedright\arraybackslash}p{(\linewidth - 2\tabcolsep) * \real{0.5000}}
  >{\raggedright\arraybackslash}p{(\linewidth - 2\tabcolsep) * \real{0.5000}}@{}}
\toprule\noalign{}
\begin{minipage}[b]{\linewidth}\raggedright
必交文件
\end{minipage} & \begin{minipage}[b]{\linewidth}\raggedright
内容
\end{minipage} \\
\midrule\noalign{}
\endhead
\bottomrule\noalign{}
\endlastfoot
\passthrough{\lstinline!code/bookstore/book\_tasks.py!} & 完成 B1--B6 后的书店程序（B2 包含两个更新方法） \\
\passthrough{\lstinline!code/iris/iris\_classification.py!} & 完成 I1--I3 后的鸢尾花程序 \\
\passthrough{\lstinline!code/bookstore/book\_support.py!} & 框架提供的公共读取、菜单、保存与显示模块 \\
\passthrough{\lstinline!code/bookstore/books.csv!} & 程序实际采用的书库输入 \\
\passthrough{\lstinline!code/bookstore/sales.csv!} & 程序实际采用的历史销售输入 \\
\passthrough{\lstinline!report.pdf!} & 两个任务的完成情况，以及鸢尾花结果和简短分析 \\
\end{longtable}

本次正式作业提交两份学生脚本，并保留运行需要的辅助模块与 CSV。\passthrough{\lstinline!practice.ipynb!}、\passthrough{\lstinline!books.txt!} 和自行创建的 \passthrough{\lstinline!hello.py!} 不用另交，验证程序 \passthrough{\lstinline!check\_book\_tasks.py!} 和 \passthrough{\lstinline!check\_iris.py!} 也不用放入压缩包。Python 环境目录、缓存和安装包不随作业提交。默认导出的 \passthrough{\lstinline!books\_updated.csv!} 无需另交；如果要采用其中的数据作为最终输入，将它另存为提交目录中的 \passthrough{\lstinline!code/bookstore/books.csv!}，并重新运行检查。具体截止时间与提交平台以课程通知为准。

\subsection{report.pdf 写什么}\label{report.pdf-ux5199ux4ec0ux4e48}

报告首页写明实验名称、学号和姓名，正文按以下三部分组织。可以从\href{report-template.md}{报告内容模板}开始，完成后导出 PDF，文件名为 \passthrough{\lstinline!report.pdf!}。

\textbf{一、完成情况。} 简要说明书店与鸢尾花两个任务是否完成。如有未完成部分或已知问题，写清具体功能；只有存在值得说明的实现取舍时才补充介绍。不必逐个粘贴书店测试用例的输出。

\textbf{二、鸢尾花结果。} 放入你生成的一张训练集散点图，保留坐标名称、单位和图例；写出本次运行的测试准确率，并以表格或等宽文本列出混淆矩阵，注明类别顺序及``行是真实类别、列是预测类别''。这些结果直接放在报告中即可，不用再交相同内容的截图集。

\textbf{三、结果分析。} 围绕自己的结果回答两点：所选两个特征下，哪些类别较容易区分，哪些仍有重叠；混淆矩阵中主要出现了哪些错分，或本次是否没有错分。每点用几句话说明即可。散点图只展示两个特征，模型使用四个特征，不能直接把图上的重叠当作模型错分的原因。

报告不要求复述讲义理论、提供安装证明或专门编写调试经历。书店功能主要通过代码运行检查，报告用于说明完成情况和展示分类结果。

\clearpage

\subsection{打包前怎样检查}\label{ux6253ux5305ux524dux600eux6837ux68c0ux67e5}

\begin{enumerate}
\def\labelenumi{\arabic{enumi}.}
\item
  在 \passthrough{\lstinline!code!} 目录运行两个项目的独立验证程序：

\begin{lstlisting}[language={}]
python bookstore/check_book_tasks.py
python iris/check_iris.py
\end{lstlisting}

  书店检查使用内置固定样例，按 B1--B6 验证对象、查询、更新、销售解析与计算，不依赖你改过的 CSV；可用 \passthrough{\lstinline!--stage B3!} 等参数单独检查。鸢尾花检查应显示 \passthrough{\lstinline!通过 3/3 项!}，核对特征选择、训练与预测调用，以及改变测试输入后的对应关系。若出现``待完成''\,``未通过''或``异常''，按提示修改对应函数后再运行。不需要提交验证截图。
\item
  运行两个作业程序，确认使用的是准备提交的文件：

\begin{lstlisting}[language={}]
python bookstore/book_tasks.py
python iris/iris_classification.py
\end{lstlisting}

  书店程序应能加载 CSV 并进入菜单；请实际尝试查询、补货、改价、保存和销售统计。鸢尾花程序应显示散点图，输出测试准确率和混淆矩阵。报告中的图和数值应来自本次提交的代码。
\item
  打开 \passthrough{\lstinline!report.pdf!}，确认图、数字和文字可读。按第 7.1 节打包后，将 ZIP 解压到另一个位置，进入其中的 \passthrough{\lstinline!code!} 目录再次运行两个程序，确认数据文件没有漏放。
\end{enumerate}

公开检查用于帮助自查，不等同于全部评分规则。若仍有未解决的问题，在报告中如实注明，保留已有实现并按课程通知提交。

\clearpage

\section{遇到问题时}\label{ux9047ux5230ux95eeux9898ux65f6}

先找到报错对应的文件或单元，再检查当前环境和输入。以下是本次实验常见的情况，Python 异常类型的进一步说明见讲义第 5 章。

\begin{longtable}[]{@{}
  >{\raggedright\arraybackslash}p{(\linewidth - 2\tabcolsep) * \real{0.5000}}
  >{\raggedright\arraybackslash}p{(\linewidth - 2\tabcolsep) * \real{0.5000}}@{}}
\toprule\noalign{}
\begin{minipage}[b]{\linewidth}\raggedright
现象
\end{minipage} & \begin{minipage}[b]{\linewidth}\raggedright
优先检查
\end{minipage} \\
\midrule\noalign{}
\endhead
\bottomrule\noalign{}
\endlastfoot
找不到 \passthrough{\lstinline!conda!} & Windows 是否打开 Miniforge / Anaconda Prompt；macOS/Linux 安装后是否重开终端或加载 conda \\
激活时提示需要 \passthrough{\lstinline!conda init!} & Windows 先使用安装器提供的 Prompt；macOS/Linux 核对第 2.4 节，不确定时提供终端类型和完整报错 \\
提示课程环境不存在 & 用 \passthrough{\lstinline!conda env list!} 核对名称，检查创建步骤是否成功 \\
下载缓慢或超时 & 按第 3.4 节区分安装器、conda 和 pip；网络稳定后重试 \\
无法导入 \passthrough{\lstinline!numpy!}、\passthrough{\lstinline!sklearn!} 或 \passthrough{\lstinline!matplotlib!} & 核对 \passthrough{\lstinline!sys.executable!}，在对应环境中用 \passthrough{\lstinline!python -m pip!} 安装；安装名为 \passthrough{\lstinline!scikit-learn!} \\
找不到 \passthrough{\lstinline!jupyterlab!} 模块 & 激活环境后安装 \passthrough{\lstinline!jupyterlab!}，再用同一个 Python 启动 \\
找不到脚本或数据文件 & 确认完整解压和当前目录；书店脚本与两份 CSV 应一起保存在 \passthrough{\lstinline!code/bookstore!} 中 \\
图未出现或程序暂停 & 确认 Matplotlib 可导入；脚本可能正在等待关闭当前图窗 \\
作业提示``待完成'' & 搜索对应 \passthrough{\lstinline!LAB01\_TODO!}；I1 选择特征列，其他位置替换占位语句，具体见项目说明 \\
函数补全后仍返回 \passthrough{\lstinline!None!} & 检查函数是否返回了所需结果；修改 Notebook 定义后重新执行相关单元 \\
改了代码但输出没变 & 检查是否保存文件、运行了另一个副本，或漏跑了 Notebook 单元 \\
符合条件的书只返回一本 & 检查是否在循环内提前 \passthrough{\lstinline!return!} \\
项目输出与默认示例不同 & 核对是否修改了 CSV；独立验证使用固定样例，不受本地数据变化影响 \\
CSV 初始导入失败 & 按提示核对表头、字段类型和书名唯一性，修正输入文件后重新运行 \\
中文读取失败或乱码 & 核对文件实际保存编码与 \passthrough{\lstinline!open!} 的编码是否一致，示例使用 UTF-8 \\
Notebook 单元一直显示运行中 & 检查是否有未结束的循环；必要时中断内核并检查代码 \\
\end{longtable}

求助时说明操作系统、运行文件、使用的终端或 Notebook，以及已经尝试的操作。安装问题提供命令和完整错误段；程序问题提供从 \passthrough{\lstinline!Traceback!} 开始到最后一行的报错。发送前遮去账号、口令和登录令牌。

\end{document}
